\documentclass[11pt,a4paper]{article}

\usepackage[margin=2.5cm]{geometry}
\usepackage{amsmath,amssymb}
\usepackage[round]{natbib}
\usepackage[colorlinks=true,linkcolor=blue,citecolor=blue,urlcolor=blue]{hyperref}

\newcommand{\Rpl}{R_{\rm pl}}
\newcommand{\Rnom}{R_{\rm nom}}
\newcommand{\Ntot}{N_{\rm L}}
\newcommand{\Nspec}{N_{\rm spec}}
\newcommand{\Trans}{\mathcal{T}}
\newcommand{\code}[1]{\texttt{#1}}

\title{Calculation of transmission spectra in \textsc{petitRADTRANS}}
\author{}
\date{\today}

\begin{document}
\maketitle

\section{Overview}

This note describes how \textsc{petitRADTRANS} (pRT) computes the transmission
spectrum, i.e.\ the wavelength-dependent transit radius $R(\nu)$ of a planet,
in the current JAX implementation
(\code{radtrans\_core/radiative\_transfer\_core.py}, functions
\code{compute\_radius\_hydrostatic\_equilibrium},
\code{compute\_transit\_radii} and
\code{compute\_transmission\_spectrum\_contribution}).
The method follows Sect.~2.2 of \citet{molliere2019} and, where possible, the
same notation is used. The calculation has four steps:
\begin{enumerate}
  \item map the pressure grid onto a radius grid assuming hydrostatic
        equilibrium (Sect.~\ref{sec:hydrostatic});
  \item compute the optical depth along grazing rays (chords) with
        impact parameter equal to the radius of each layer
        (Sect.~\ref{sec:tau});
  \item compute the band-averaged transmission of each chord, using the
        correlated-$k$ treatment if applicable (Sect.~\ref{sec:transmission});
  \item integrate the opaque fraction of each annulus to obtain an
        effective planetary area and hence the transit radius
        (Sect.~\ref{sec:area}).
\end{enumerate}
The transmission contribution function is described in
Sect.~\ref{sec:contribution}. All steps are performed independently for each
frequency $\nu$ (or correlated-$k$ bin); for clarity the frequency index is
suppressed below.

\paragraph{Grid and indexing.}
The atmosphere is described on $\Ntot$ layers with pressures
$P_1 < P_2 < \dots < P_{\Ntot}$, so that $i=1$ is the top of the atmosphere
and $i=\Ntot$ the bottom. Temperatures $T_i$ and mean molar masses $\mu_i$
are given at each layer. Following \citet{molliere2019}, $r_i$ denotes the
radius of layer $i$ measured from the planetary centre, $\kappa$ denotes a
mass opacity (cm$^2$\,g$^{-1}$) and $\rho_i$ the mass density,
\begin{equation}
  \rho_i = \frac{P_i\,\mu_i\,m_{\rm u}}{k_{\rm B}\,T_i},
  \label{eq:rho}
\end{equation}
with $m_{\rm u}$ the atomic mass unit and $k_{\rm B}$ the Boltzmann constant
(cgs units throughout).

\section{Hydrostatic radius grid}
\label{sec:hydrostatic}

The radius grid is obtained by integrating the hydrostatic equilibrium
equation, $\mathrm{d}P/\mathrm{d}r = -\rho\, g(r)$, from a reference point
$r(P_0) = \Rpl$ at which the gravity is $g_0 \equiv g(\Rpl)$.

\paragraph{Variable gravity} (default, $g(r) = g_0\Rpl^2/r^2$). Here
$\mathrm{d}(1/r) = \mathrm{d}P/(\rho\, g_0 \Rpl^2)$, so that
\begin{equation}
  \frac{1}{r_i} = \frac{1}{\Rpl}
  + \frac{1}{g_0\Rpl^2}\int_{P_0}^{P_i}\frac{\mathrm{d}P}{\rho} .
  \label{eq:hydro_var}
\end{equation}

\paragraph{Constant gravity} ($g(r) = g_0$):
\begin{equation}
  r_i = \Rpl - \frac{1}{g_0}\int_{P_0}^{P_i}\frac{\mathrm{d}P}{\rho} .
  \label{eq:hydro_const}
\end{equation}

\paragraph{Discretisation.}
The integral of $1/\rho$ over each pressure interval is computed by
piecewise-parabolic quadrature. For $i = 1,\dots,\Ntot-2$,
\begin{equation}
  I_i \equiv \int_{P_i}^{P_{i+1}} \frac{\mathrm{d}P}{\rho}
  \approx \int_{P_i}^{P_{i+1}} p_i(P)\,\mathrm{d}P ,
\end{equation}
where $p_i(P)$ is the quadratic passing through
$(P_i,1/\rho_i)$, $(P_{i+1},1/\rho_{i+1})$ and $(P_{i+2},1/\rho_{i+2})$.
The last interval uses the trapezoidal rule,
$I_{\Ntot-1} = \tfrac12\,(1/\rho_{\Ntot-1} + 1/\rho_{\Ntot})(P_{\Ntot}-P_{\Ntot-1})$.
The integral from the bottom of the atmosphere to layer $i$ is then the
cumulative sum $\sum_{k=i}^{\Ntot-1} I_k$.

The offset to the reference pressure is obtained by locating the interval
$k$ with $P_k \le P_0 < P_{k+1}$ and integrating the same parabola $p_k$
(or, if $k = \Ntot-1$, the trapezoid with end-point values
$1/\rho_{\Ntot-1}$ and $1/\rho_{\Ntot}$) from $P_0$ to $P_{k+1}$. This
fixes the integration constant such that Eqs.~\eqref{eq:hydro_var}
and~\eqref{eq:hydro_const} are satisfied. If $P_0$ does not lie inside the
pressure grid, the offset is set to zero, which places $\Rpl$ at the bottom
layer, $r_{\Ntot} = \Rpl$.

In the variable-gravity case, Eq.~\eqref{eq:hydro_var} yields $1/r_i \le 0$
for the uppermost layers if the atmosphere is very extended (low gravity,
high temperature). In \code{compute\_transit\_radii}, any negative $r_i$ is
replaced by the radius of the uppermost layer with $r_i \ge 0$, i.e.\ the
affected layers are collapsed onto that radius and no longer contribute to
the spectrum.

\section{Optical depth along grazing rays}
\label{sec:tau}

Consider a ray passing the planet at impact parameter $r_i$. Between two
consecutive layers $j$ and $j+1$ (with $j < i$), the ray travels a distance
on each side of the point of closest approach of
\begin{equation}
  \Delta s_{ij} = \sqrt{r_j^2 - r_i^2} - \sqrt{r_{j+1}^2 - r_i^2} .
  \label{eq:ds}
\end{equation}
Approximating the extinction coefficient in this shell by the trapezoidal
mean $(\kappa_j\rho_j + \kappa_{j+1}\rho_{j+1})/2$, and multiplying by $2$ to
account for the ingoing and outgoing halves of the chord, gives Eq.~(6) of
\citet{molliere2019},
\begin{equation}
  \tau_i = \sum_{j=1}^{i-1}
  \left(\kappa_j\rho_j + \kappa_{j+1}\rho_{j+1}\right)\Delta s_{ij},
  \qquad i = 2,\dots,\Ntot,
  \label{eq:tau}
\end{equation}
with $\tau_1 = 0$ for the ray grazing the top of the atmosphere. In the
code, Eq.~\eqref{eq:tau} is evaluated as a single matrix product of the
$(\Ntot-1)\times(\Ntot-1)$ lower-triangular path-length matrix
$\Delta s_{ij}$ with the vector of paired extinctions
$\kappa_j\rho_j + \kappa_{j+1}\rho_{j+1}$, separately for each frequency,
$g$-point and species.

\paragraph{Opacity entering $\tau$.}
Eq.~\eqref{eq:tau} is evaluated for every species $s$ and correlated-$k$
point $l$, giving $\tau_{i,s,l}$. The opacity $\kappa_{j,s,l}$ is the
mass-fraction-weighted line opacity of species $s$, $X_s\,\kappa^{\rm line}_{j,s,l}$.
The continuum is added to the first species slot ($s=1$) only:
\begin{equation}
  \kappa_{j,1,l} = X_1\,\kappa^{\rm line}_{j,1,l}
  + \kappa^{\rm abs}_{{\rm cont},j} + \kappa^{\rm scat}_{{\rm cont},j},
\end{equation}
where $\kappa^{\rm abs}_{\rm cont}$ contains all continuum absorption (e.g.\
collision-induced absorption, H$^-$, cloud absorption, grey cloud decks)
and $\kappa^{\rm scat}_{\rm cont}$ all continuum scattering (e.g.\ Rayleigh
and cloud scattering). In transmission, scattering is treated as pure
extinction: scattered photons are assumed to be removed from the beam.
Since the continuum is grey within a correlated-$k$ bin, adding it to a
single species is equivalent to multiplying the total transmission by
$\exp(-\tau^{\rm cont}_i)$, provided $\sum_l \Delta g_l = 1$.

If the line opacities have already been combined into a single effective
opacity upstream (summed line-by-line opacities, or correlated-$k$
opacities combined with the random-overlap method), only the first species
slot is used, i.e.\ effectively $\Nspec = 1$.

\section{Transmission of a chord}
\label{sec:transmission}

The transmission of the chord with impact parameter $r_i$, averaged over the
correlated-$k$ bin, is computed as in Eq.~(2) of \citet{molliere2019},
assuming that the opacity distributions of different species are
uncorrelated:
\begin{equation}
  \bar{\Trans}_i = \prod_{s=1}^{\Nspec}
  \left[\sum_{l=1}^{N_g} \Delta g_l\, e^{-\tau_{i,s,l}}\right]
  = \prod_{s=1}^{\Nspec} \bar{\Trans}_{i,s} ,
  \label{eq:transm}
\end{equation}
where $\Delta g_l$ are the Gaussian quadrature weights in $g$-space. In
line-by-line mode $N_g = 1$ and the $g$-sum is omitted, so that
$\bar{\Trans}_i = \exp(-\sum_s\tau_{i,s})$.

\section{Effective area and transit radius}
\label{sec:area}

The effective area of the planet is given by Eq.~(7) of
\citet{molliere2019},
\begin{equation}
  A = 2\pi\int_0^{\Rpl} r\,\left(1 - \Trans(r)\right)\mathrm{d}r .
  \label{eq:area}
\end{equation}
The planet is assumed to be opaque below the lowest layer
($\Trans = 0$ for $r < r_{\Ntot}$), which contributes $\pi r_{\Ntot}^2$.
The remaining integral over the atmosphere, $r_{\Ntot} \le r \le r_1$, is
evaluated with the trapezoidal rule on the layer radii:
\begin{equation}
  A = \pi r_{\Ntot}^2
  + \pi\sum_{i=1}^{\Ntot-1}
  \left[r_i\left(1-\bar{\Trans}_i\right)
      + r_{i+1}\left(1-\bar{\Trans}_{i+1}\right)\right]
  \left(r_i - r_{i+1}\right).
  \label{eq:area_disc}
\end{equation}
Here the factor $2$ of Eq.~\eqref{eq:area} has cancelled against the factor
$1/2$ of the trapezoidal rule. The transit radius is the radius of an opaque
disc of the same area,
\begin{equation}
  R = \sqrt{\frac{A}{\pi}} =
  \left\{ r_{\Ntot}^2 + \sum_{i=1}^{\Ntot-1}
  \left[r_i\left(1-\bar{\Trans}_i\right)
      + r_{i+1}\left(1-\bar{\Trans}_{i+1}\right)\right]
  \left(r_i - r_{i+1}\right) \right\}^{1/2} ,
  \label{eq:radius}
\end{equation}
which is what \code{compute\_transit\_radii} returns, together with the
radius grid $r_i$. Because $\bar{\Trans}_1 = 1$ ($\tau_1 = 0$), the topmost
layer contributes only through the trapezoid of the first annulus.

\section{Transmission contribution function}
\label{sec:contribution}

Following Eq.~(8) of \citet{molliere2019}, the contribution of layer $i$ is
estimated by removing its opacity and recomputing the transit radius:
\begin{equation}
  C_{\rm tr}^{i} =
  \frac{\Rnom^2 - R^2(\kappa_i = 0)}
       {\sum_{j=1}^{\Ntot}\left[\Rnom^2 - R^2(\kappa_j = 0)\right]} ,
  \label{eq:contrib}
\end{equation}
where $\Rnom$ is the nominal transit radius and $R(\kappa_i=0)$ the radius
obtained when the opacity of layer $i$ (all species, all $g$-points and the
continuum) is set to zero. The radius grid $r_i$ is kept fixed. Squared
radii are used since the transit depth scales with the planet's area. If
the denominator vanishes, $C_{\rm tr}^i$ is set to zero.

Rather than recomputing Eq.~\eqref{eq:tau} $\Ntot$ times, the code uses the
fact that layer $i$ only enters the shells above and below it, $j = i-1$
and $j = i$. Removing its opacity therefore changes the optical depth of
every chord $k$ by a rank-one correction,
\begin{equation}
  \tau_k(\kappa_i = 0) = \tau_k
  - \kappa_i\rho_i\left(\Delta s_{k,i-1} + \Delta s_{k,i}\right),
  \label{eq:tau_loo}
\end{equation}
with $\Delta s_{k,j} \equiv 0$ outside $1 \le j < k$. The modified optical
depths are then passed through Eqs.~\eqref{eq:transm}
and~\eqref{eq:radius}. Here the continuum scattering optical depth is
applied as a separate factor $\exp(-\tau^{\rm scat}_k)$ on the total
transmission, which is equivalent to the treatment in
Sect.~\ref{sec:tau}. The nominal radius $\Rnom$ in Eq.~\eqref{eq:contrib} is
recomputed inside the same function with the same numerics, rather than
taken from the main transmission calculation, so that differences in
numerical precision between the two do not appear as a spurious,
layer-independent contribution.

\section{Numerical notes}

\begin{itemize}
  \item The radius grid is computed in the precision of the input pressure
        grid (double precision when JAX's 64-bit mode is enabled). The
        optical depths, transmissions and the area integral are evaluated in
        a reduced working precision (\code{WORKING\_JAX\_DTYPE}, currently
        single precision) to save memory and time.
  \item The square roots in Eq.~\eqref{eq:ds} are evaluated as
        $\sqrt{\max(x,0)}$ with a gradient-safe branch, so that automatic
        differentiation does not produce NaNs for $r_j \le r_i$.
  \item The path-length matrix is explicitly set to zero for $j \ge i$.
        Although $\sqrt{r_i^2 - r_i^2}$ should vanish, the compiled
        expression can return a small rounding residual of $r_i^2$, whose
        square root is a spurious path of order $1\,$cm. For very opaque
        layers (e.g.\ a grey cloud deck) this would otherwise produce an
        optical depth $\tau \sim 10^{14}$ in a chord that should not cross
        the layer.
\end{itemize}

\begin{thebibliography}{1}
\bibitem[Molli\`ere et al.(2019)]{molliere2019}
P.~Molli\`ere, J.~P. Wardenier, R.~van Boekel, Th.~Henning,
K.~Molaverdikhani, and I.~A.~G. Snellen,
\newblock petitRADTRANS: a Python radiative transfer package for exoplanet
characterization and retrieval,
\newblock \emph{A\&A}, 627, A67 (2019),
\newblock \href{https://doi.org/10.1051/0004-6361/201935470}{doi:10.1051/0004-6361/201935470}.
\end{thebibliography}

\end{document}
